\documentclass[aspectratio=169]{beamer}

% --- Temas y Colores ---
\usetheme{Madrid} % Tema clásico y formal para investigación
\usecolortheme{default} % Puedes cambiarlo a 'whale', 'dolphin', etc.

% --- Paquetes y Librerías Necesarias ---
\usepackage[utf8]{inputenc} % Codificación para tildes y eñes
\usepackage[spanish, es-tabla]{babel} % Idioma español
\usepackage{graphicx} % Para insertar imágenes
\usepackage{amsmath, amssymb, amsfonts} % Para fórmulas matemáticas (FCME, Gram-Schmidt)
\usepackage{booktabs} % Para tablas con formato profesional (toprule, midrule, bottomrule)
\usepackage{tikz} % Excelente para crear diagramas de arquitectura en código
\usepackage{hyperref} % Para hipervínculos
\usepackage{listings} % Por si necesitas mostrar fragmentos de código en C99 o Python

% --- Configuración de la Portada ---
\title[Sensado SDR Autónomo - ANE]{Desarrollo y Despliegue de un Sistema de Sensado SDR Autónomo de Alto Rendimiento}
\subtitle{Monitoreo Institucional del Espectro Radioeléctrico}
\author[Tu Nombre]{Tu Nombre / Tu Equipo}
\institute[Tu Institución]{Nombre de tu Grupo de Investigación o Universidad \\ \vspace{0.2cm} Proyecto en contexto con la ANE}
\date{\today}

\begin{document}

% --- Diapositiva: Portada ---
\begin{frame}
    \titlepage
\end{frame}

% --- Diapositiva: Índice (Opcional, pero recomendado) ---
\begin{frame}{Contenido}
    \tableofcontents
\end{frame}

% --- Diapositiva 1: Motivación ---
\section{Motivación}
\begin{frame}{Motivación: Contexto Global}
    \begin{itemize}
        \item \textbf{El Recurso:} El espectro radioeléctrico es un bien escaso y estratégico.
        \item \textbf{Rol de la ANE:} Garantizar el uso legal y eficiente (TDT, bandas móviles, FM).
        \item \textbf{Necesidad Global:} Demanda por sensores autónomos, de bajo costo y alta disponibilidad para redes nacionales masivas.
    \end{itemize}
\end{frame}

% --- Diapositiva 2: Problema ---
\section{Problema}
\begin{frame}{El Problema: La Brecha Técnica}
    \begin{itemize}
        \item \textbf{Carga Computacional:} Procesar flujos IQ de alta velocidad (40 MB/s) en hardware embebido (Raspberry Pi 5) sin \textit{sample drops}.
        \item \textbf{Artefactos de Hardware:} Receptores SDR económicos presentan fallas intrínsecas (DC Spike, desbalance IQ) que contaminan mediciones.
        \item \textbf{Entornos Inestables:} Operación remota con enlaces LTE/4G precarios y cortes abruptos de energía.
    \end{itemize}
\end{frame}

% --- Diapositiva 3: Estado del Arte ---
\section{Estado del Arte}
\begin{frame}{Estado del Arte: SDR de Bajo Costo y Edge Computing}
    \textbf{Desarrollos Actuales (SOTA):}
    \begin{itemize}
        \item Redes Crowdsourced (ej. Electrosense con RTL-SDR y RPi).
        \item Nodos para bandas puntuales (TVS, TETRA) entre \$150 - \$350 USD.
    \end{itemize}
    \vspace{0.3cm}
    \textbf{Limitaciones para la ANE:}
    \begin{itemize}
        \item \textit{Cuellos de Botella:} GNU Radio / Python saturan la CPU en el borde.
        \item \textit{Falta de Rigor Metrológico:} Sin limpieza espectral en tiempo real.
        \item \textit{Ausencia de Grado Institucional:} Sin estrategias de despliegue seguro o persistencia.
    \end{itemize}
\end{frame}

% --- Diapositiva 4: Objetivos ---
\section{Objetivos}
\begin{frame}{Objetivos del Proyecto}
    \textbf{Meta General:} Desarrollar un sistema de sensado SDR autónomo de alto rendimiento para el monitoreo institucional. \\
    \vspace{0.3cm}
    \textbf{Hitos Específicos:}
    \begin{enumerate}
        \item Implementar una arquitectura híbrida optimizada (C99/Python).
        \item Garantizar la integridad de datos (escritura atómica, colas FIFO).
        \item Validar metrológicamente el sistema con estándares de laboratorio.
        \item Automatizar la documentación técnica para mantenimiento.
    \end{enumerate}
\end{frame}

% --- Diapositiva 5: Metodología I ---
\section{Metodología}
\begin{frame}{Metodología I: Arquitectura y Stack Tecnológico}
    % Aquí puedes insertar tu imagen o usar TikZ
    \begin{center}
        \includegraphics[width=0.6\textwidth]{example-image} % Reemplaza 'example-image' con el nombre de tu diagrama
    \end{center}
    \begin{itemize}
        \item \textbf{Hardware:} Raspberry Pi 5, HackRF One, Módulo SIM7600G-H.
        \item \textbf{Software (Híbrido):} Python 3 (Control/Orquestación) + C99 y FFTW3 (Plano de Datos DSP).
        \item \textbf{Comunicación:} ZeroMQ (ZMQ) para latencia ultra baja.
    \end{itemize}
\end{frame}

% --- Diapositiva 6: Metodología II ---
\begin{frame}{Metodología II: Despliegue y Algoritmos}
    \textbf{Despliegue e Infraestructura:}
    \begin{itemize}
        \item OTA-ANE v7 con VPN WireGuard y Port Knocking.
    \end{itemize}
    \vspace{0.3cm}
    \textbf{Algoritmos Clave (DSP en el Borde):}
    \begin{itemize}
        \item \textit{Limpieza Espectral:} Offset \& Crop y Spectral Stitching (eliminación del DC Spike).
        \item \textit{Cumplimiento Normativo:} Algoritmo FCME para estimación del piso de ruido.
        \item \textit{Metrología:} Ortogonalización de Gram-Schmidt para desbalance IQ.
    \end{itemize}
\end{frame}

% --- Diapositiva 7: Resultados I ---
\section{Resultados y Discusión}
\begin{frame}{Resultados: Evidencia Metrológica}
    \begin{itemize}
        \item \textbf{Piso de Ruido (DANL):} Caracterización estadística ($N=10,000$) validada con pruebas ANOVA entre unidades SDR.
        \item \textbf{Precisión de Potencia:} Error promedio $< 1.5$ dB (calibración Keysight).
        \item \textbf{Exactitud de Frecuencia:} Errores $< 13$ ppm (sincronización vía tono GSM FCCH).
    \end{itemize}
    \vspace{0.3cm}
    % Ejemplo de tabla de resultados
    \begin{table}
        \centering
        \begin{tabular}{lcc}
            \toprule
            \textbf{Parámetro} & \textbf{SOTA (Promedio)} & \textbf{Nuestro Sistema} \\
            \midrule
            RMSE Potencia & $> 3$ dB & $< 1$ dB \\
            Pérdida de Muestras & Alta & $0\%$ a 40 MB/s \\
            \bottomrule
        \end{tabular}
        \caption{Comparativa de Rendimiento}
    \end{table}
\end{frame}

% --- Diapositiva 8: Resultados II ---
\begin{frame}{Resultados y Discusión: Operatividad}
    \begin{itemize}
        \item \textbf{Evidencia de Software:} 168 pruebas automatizadas exitosas (tolerancia a fallos, concurrencia).
        \item \textbf{Streaming:} Distribución de audio demodulado (AM/FM) de baja latencia vía WebRTC y Opus.
        \item \textbf{Automatización:} Generación de reportes JSON de cumplimiento (elimina subjetividad).
    \end{itemize}
\end{frame}

% --- Diapositiva 9: Conclusiones ---
\section{Conclusiones}
\begin{frame}{Conclusiones}
    \begin{itemize}
        \item \textbf{Confirmación Técnica:} Sensor de alto rendimiento que supera los cuellos de botella del SOTA y compite con instrumentos de laboratorio en linealidad (RMSE $< 1$ dB).
        \item \textbf{Impacto Real:} Descentralización efectiva del monitoreo espectral, permitiendo a la ANE gestionar mediciones masivas, remotas y a una fracción del costo.
    \end{itemize}
\end{frame}

% --- Diapositiva 10: Trabajo Futuro ---
\section{Trabajo Futuro}
\begin{frame}{Trabajo Futuro}
    \begin{itemize}
        \item Integración de \textbf{GPSDO} (Osciladores Disciplinados por GPS) para estabilidad de frecuencia atómica.
        \item Implementación de \textbf{pruebas unitarias nativas} para el motor DSP en C.
        \item Exploración de \textbf{aceleración por GPU} (Edge AI) para cálculos de PSD a gran escala.
    \end{itemize}
\end{frame}

\end{document}